%% ma: L10-B13-0060
%% lop: 10
%% chuong: 5
%% bai: 13
%% dang: 10.13.7
%% loai: TLN
%% muc: TH
%% dap_an: "5"
%% nguon: "chuyen-de-cac-so-dac-trung-cua-mau-so-lieu-khong-ghep-nhom-toan-10-chuong-trinh-moi.pdf (GV. Phan Nhật Linh) (P02), Bài 02 Dạng 1 Phần III Câu 10"
%% kien_thuc: "Lớp 10 Bài 13 (đọc và phân tích số liệu), tỉ lệ phần trăm"
%% trang_thai: chinh-thuc
%% ghi_chu: "Tự duyệt theo yêu cầu giáo viên 22/10/2026."
\begin{ex}
Trong $6$ tháng đầu năm, số sản phẩm bán ra mỗi tháng của một cửa hàng đều tăng khoảng $25\%$ so với tháng trước đó. Biết rằng trong bảng dưới đây, số sản phẩm bán ra của đúng một tháng bị nhập sai. Tháng đó là tháng thứ mấy?
\begin{center}
\begin{tabular}{|c|c|c|c|c|c|c|}\hline
Tháng & $1$ & $2$ & $3$ & $4$ & $5$ & $6$\\\hline
Số sản phẩm & $145$ & $180$ & $225$ & $279$ & $390$ & $435$\\\hline
\end{tabular}
\end{center}
\shortans{5}
\loigiai{Tỉ lệ tăng so với tháng trước: $24{,}1\%$; $25\%$; $24\%$; $39{,}8\%$; $11{,}5\%$.

Tháng 5 và tháng 6 đều lệch xa $25\%$, mà $279\cdot1{,}25\approx349$ và $349\cdot1{,}25\approx436\approx435$, nên số liệu tháng $5$ bị nhập sai.}
\end{ex}
